% ---------------------------------------------------------------------
%  Title page
% ---------------------------------------------------------------------
\begin{titlepage}
\thispagestyle{empty}
\begin{center}
\vspace*{0.6cm}
{\sffamily\color{agrey}\small A TEACHING TEXT\par}
\vspace{0.5cm}
{\sffamily\bfseries\fontsize{34}{38}\selectfont\color{ablue} Amber and\\[2pt] 6G Ambient IoT\par}
\vspace{0.5cm}
{\sffamily\Large From backscatter physics and energy harvesting\\
to protocols and a discrete-event simulator\par}

\vspace{1.0cm}
\begin{tikzpicture}[scale=1.2,every node/.style={font=\sffamily\scriptsize}]
  % ground line
  \draw[agrey!50] (-0.8,0) -- (9.6,0);
  % base-station mast
  \draw[line width=1.6pt,ablue] (0,0) -- (0,2.1);
  \fill[ablue] (-0.22,2.1) rectangle (0.22,2.65);
  \node[below] at (0,0) {base station};
  % carrier arcs towards the tag
  \foreach \r/\o in {0.9/0.95,1.5/0.75,2.1/0.55,2.7/0.38}{
    \draw[ablue,opacity=\o,line width=1.2pt]
      ([shift={(-24:\r)}]0.25,2.38) arc (-24:24:\r);}
  % passive node on a post
  \draw[line width=1pt,agrey] (4.7,0) -- (4.7,1.45);
  \draw[line width=1.1pt,fill=lightgreen] (4.25,1.45) rectangle (5.15,2.15);
  \draw[line width=1.1pt] (4.7,2.15) -- (4.7,2.75);
  \draw[line width=1.1pt] (4.45,2.75) -- (4.95,2.75);
  \node[below] at (4.7,0) {passive node};
  % backscatter arcs going back
  \foreach \r/\o in {0.5/0.95,0.95/0.7,1.4/0.45}{
    \draw[aorange,opacity=\o,line width=1.2pt]
      ([shift={(156:\r)}]4.55,2.5) arc (156:204:\r);}
  % capacitor saw-tooth
  \draw[->,agrey] (6.4,0.2) -- (6.4,2.5) node[above] {$V_C$};
  \draw[->,agrey] (6.4,0.2) -- (9.4,0.2) node[right] {$t$};
  \draw[ablue,line width=1.1pt] (6.4,0.2) .. controls (6.8,1.4) .. (7.2,1.95)
     -- (7.3,1.1) .. controls (7.7,1.7) .. (8.0,1.95)
     -- (8.1,1.1) .. controls (8.5,1.7) .. (8.8,1.95)
     -- (8.9,1.1);
  \draw[dashed,agrey] (6.4,1.95) -- (9.4,1.95) node[right] {$V_{\mathrm{high}}$};
  \draw[dashed,agrey] (6.4,1.1) -- (9.4,1.1) node[right] {$V_{\mathrm{low}}$};
\end{tikzpicture}

\vfill
{\sffamily\large
A study guide to the SenSys 2026 demo abstract\\[3pt]
\emph{``Demo Abstract: Amber: 6G Ambient IoT Simulator''}\\[2pt]
M.\ Manafova, T.\ Voigt and L.\ Mottola\\[10pt]
and to the open-source Amber simulator, version~0.1.0\par}
\vspace{0.6cm}
{\sffamily\small\color{agrey}
This is an independent teaching text. It was not written by the Amber authors.\\
Statements about the code were checked against the source; the behaviours that matter were also tested by running it.\par}
\vspace{0.4cm}
\end{center}
\end{titlepage}

\frontmatter

% ---------------------------------------------------------------------
\chapter*{Preface}
\addcontentsline{toc}{chapter}{Preface}
\markboth{Preface}{Preface}

\paragraph{The problem this book solves.}
A device that has no battery, harvests a few microwatts from radio waves, and
answers a base station by \emph{reflecting} the base station's own signal is
easy to describe in one sentence and surprisingly hard to reason about. Its
ability to talk depends on its stored energy; its stored energy depends on
the radio channel; the radio channel depends on distance, antennas and on
\emph{what the device itself is doing} at that moment; and the protocol that
sits on top must cope with hundreds of such devices that wake up and fall
asleep at unpredictable times. No single textbook chapter covers all of this.
The \emph{Amber} simulator was built to let researchers explore exactly this
tangle, and its authors describe it in a two-page SenSys demo abstract.

\paragraph{What this book does.}
We use Amber as a vehicle for teaching the \emph{whole} subject. For every
part of the simulator we first explain the \emph{theory} it abstracts---the
physics of antennas, path loss, backscatter, rectifiers and capacitors, the
mathematics of random access and interference cancellation, the machinery of
discrete-event simulation---and only then show \emph{how Amber models it},
where the model is exact, where it is a deliberate simplification, and where
the code does something different from what the textbook says. You will not
find page after page of source listings: the aim is that after reading you
could re-derive every equation inside Amber and write a simulator like it
yourself.

\paragraph{How the book is organised.}
\begin{description}[leftmargin=1.2cm,style=nextline,font=\sffamily\bfseries\color{ablue}]
 \item[Part I (Chapters 1--2).] The idea of Ambient IoT, 3GPP's view of it,
   and a one-chapter tour of Amber so that you always know where a topic fits.
 \item[Part II (Chapters 3--8).] The physics: decibels and noise, antennas,
   path loss, backscatter communication, RF energy harvesting, and energy
   storage in a capacitor with threshold-driven intermittent operation.
 \item[Part III (Chapters 9--11).] Sharing the medium: ALOHA and its framed
   variants, capture, interference and successive interference cancellation,
   and how these combine into registration and polling protocols.
 \item[Part IV (Chapters 12--13).] Simulation methodology (discrete-event
   simulation with SimPy) and the architecture of Amber.
 \item[Part V (Chapters 14--18).] A module-by-module deep dive: topology and
   coverage, the node controller, the backscatter module, the base-station
   engine, and packet analysis.
 \item[Part VI (Chapters 19--20).] Experiments: the four shipped demos, our
   re-creation of the paper's node-density experiment, and how to build your
   own protocol on top of Amber.
 \item[Part VII (Chapter 21).] A critical assessment: assumptions, limits and
   a catalogue of subtle behaviours that anyone using the code should know.
 \item[Appendices.] Derivations, a formula and parameter sheet, an annotated
   reading of the paper, exercise solutions and a glossary.
\end{description}

\paragraph{Conventions used throughout.}
The coloured boxes have fixed meanings.
\begin{keyidea}[Key idea]
A concept you must carry forward. If you only read the boxes, you still get the skeleton of the subject.
\end{keyidea}
\begin{intuition}
A physical or mathematical way of \emph{seeing} why a result is true, usually before or after a derivation.
\end{intuition}
\begin{example}[Worked example]
A calculation with real Amber numbers, each checked against the simulator or by an independent hand calculation.
\end{example}
\begin{reality}[Code reality check]
Where the code deviates from the textbook theory, or from what its own documentation says. These are the most valuable boxes if you intend to \emph{use} Amber for research. Each was reproduced by a small experiment (see Chapter~\ref{ch:limits}).
\end{reality}
\begin{codemap}
A short table linking a concept to the module, class or function that implements it.
\end{codemap}
Identifiers set in \code{typewriter} are names in the Amber source. Section
and chapter exercises are collected at the end of each chapter; worked
solutions are in Appendix~\ref{app:solutions}.

\paragraph{A note on honesty.}
Amber is a young research prototype (version~0.1.0). The accompanying paper is
a demo abstract: it deliberately describes the simulator at a high level and
reports one experiment. We therefore draw a clear line between (i) what the
paper states, (ii) what the code implements, and (iii) what we ran ourselves
to find out. Numbers reported as \emph{measured} come from runs of the
unmodified simulator; their scripts are in the \file{textbook/code}
directory, so you can reproduce every figure and table.

\paragraph{Prerequisites.}
Undergraduate-level probability (random variables, the Poisson distribution),
basic circuit analysis (Ohm's law, a capacitor's $i=C\,\mathrm{d}v/\mathrm{d}t$),
elementary Python, and a first course in communications or signals. We
introduce everything else.

% ---------------------------------------------------------------------
\chapter*{Reading Roadmap}
\addcontentsline{toc}{chapter}{Reading roadmap}
\markboth{Reading roadmap}{Reading roadmap}

Different readers need different routes. The table suggests three.

\begin{center}
\small
\begin{tabularx}{\linewidth}{@{}>{\sffamily\bfseries}p{2.9cm}X@{}}
\toprule
If you are\dots & Read\dots \\
\midrule
New to Ambient IoT and want the concepts &
 Chapters 1, 2, 3, 6, 7, 8, then 9--11. Skim Part V.\\[3pt]
A wireless/networking researcher who will \emph{use} Amber &
 Chapters 2, 6, 8, 10, 12, 13, then all of Part V and Chapter 21.
 Do Chapter 19 and 20 hands-on.\\[3pt]
Extending Amber (new protocol, energy model, channel) &
 Chapters 13, 15--17, 20, 21, and Appendix~\ref{app:ref}. The
 \emph{Code reality check} boxes in Chapters 8, 10, 14 and 17 matter most.\\
\bottomrule
\end{tabularx}
\end{center}

\begin{keyidea}[The one-paragraph summary of the whole book]
An Ambient IoT node lives a double life. In the \emph{energy} domain it is a
capacitor that a weak radio-frequency (or solar) source fills slowly and a
duty-cycled load drains quickly, switching on at $V_{\mathrm{high}}$ and off at
$V_{\mathrm{low}}$. In the \emph{communication} domain it is a passive reflector that can answer a
base station only when it is awake, only in the slot the base station
advertised, and only with the power budget the two-way path loss leaves it.
The base station cannot see the node's energy state, so the protocol must be
robust to nodes that are absent, collide, or arrive late. Amber simulates the
coupling between these two domains millisecond by millisecond.
\end{keyidea}

% ---------------------------------------------------------------------
\chapter*{Notation and Symbols}
\addcontentsline{toc}{chapter}{Notation and symbols}
\markboth{Notation}{Notation}

Quantities are written in the SI system unless stated. Logarithmic powers use
\si{\dBm} (decibels relative to one milliwatt), antenna gains use \si{\dBi}
(relative to an isotropic radiator), and pure ratios use \si{\dB}. The right
column shows the corresponding name in the Amber source where there is one.

\begin{center}
\small
\renewcommand{\arraystretch}{1.12}
\begin{tabularx}{\linewidth}{@{}>{$}l<{$}Xl@{}}
\toprule
\text{Symbol} & Meaning & Amber name\\
\midrule
f_c,\ \lambda,\ c & carrier frequency, wavelength $\lambda=c/f_c$, speed of light & \code{freq_hz}, \code{C}\\
d,\ d_{2D},\ d_{3D} & distance; horizontal and slant (3-D) distance & \code{distance_m}\\
h_{BS},\ h_{UT} & antenna heights of base station and node & \codew{Sector.height}, \code{Node.height}\\
G,\ G_{\max} & antenna gain (linear); peak gain & \codew{antenna_gain_dbi}\\
\varphi,\ \varphi_{3\mathrm{dB}},\ A_m & angular error, half-power beamwidth, pattern floor & \codew{beamwidth_deg}\\
L(d),\ \mathrm{PL}(d) & path loss (linear ratio $\le1$; dB) & \code{fspl_lin}, \codew{urban_macro_loss_db_38_901}\\
P_{BS} & transmit power of a base-station sector & \codew{Sector.power}\\
P_{\mathrm{dl}} & downlink power arriving at a node & \code{best_pw_dbm}\\
\eta & power split: fraction reflected (backscattered) & \codew{Node.efficiency}\\
P_h,\ P_{bs} & harvested and backscattered power at a node & \codew{harvesting_power_dbm}, \codew{backscatter_power_dbm}\\
P_{\mathrm{ul}} & backscatter power received at the base station & \codew{per_sector_powers}\\
k_B,\ T_0,\ B,\ F & Boltzmann constant, 290\,K, noise bandwidth, noise factor & \codew{thermal_noise_watts}\\
N,\ \gamma,\ \gamma_{\mathrm{req}} & noise power, SINR, required SINR & \codew{required_sinr_db}\\
\varepsilon & SIC residual: $1-$ cancellation factor & \codew{cancellation_factor}\\
C,\ V,\ E=\tfrac12CV^2 & capacitance, capacitor voltage, stored energy & \codew{CapacitorParams.C}\\
\Rs,\ \Rl & series and leakage resistance & \code{R_series}, \code{R_leakage}\\
\Vs & equivalent source voltage of the harvester & \codew{voltage_source}\\
\Vhigh,\ \Vlow,\ \Vmax & activation, shut-down and clamp voltages & \codew{thresholds_v}, \code{voltage_max}\\
\tau_c=\Rs C,\ \tau_\ell=\Rl C & charging and leakage time constants & --\\
\Delta t & simulation time step (1\,ms) & \code{dt}\\
I_L & load current of the node & \code{CurrentsA}\\
S,\ n,\ K & slots per frame, contending nodes, number of subcarrier shifts & \code{rx_slots}, \codew{subcarrier_shift}\\
g & offered load of a random-access system (attempts per slot) & --\\
\bottomrule
\end{tabularx}
\end{center}
